\documentclass[runningheads]{llncs}
%
\usepackage[T1]{fontenc}
% T1 fonts will be used to generate the final print and online PDFs,
% so please use T1 fonts in your manuscript whenever possible.
% Other font encondings may result in incorrect characters.
%
\usepackage{graphicx}
% Used for displaying a sample figure. If possible, figure files should
% be included in EPS format.
%
\usepackage{amsmath,amssymb}
\usepackage{algorithm}
\usepackage{algorithmic}
\usepackage{booktabs}
\usepackage{multirow}
%
% If you use the hyperref package, please uncomment the following two lines
% to display URLs in blue roman font according to Springer's eBook style:
%\usepackage{color}
%\renewcommand\UrlFont{\color{blue}\rmfamily}
%\urlstyle{rm}
%
\usepackage{hyperref}

\begin{document}
%
\title{AMT-UNet: Adaptive Masked Transformer U-Net for Multi-Task Brain Tumor Segmentation and Classification}
%
\titlerunning{AMT-UNet for Brain Tumor Analysis}
% If the paper title is too long for the running head, you can set
% an abbreviated paper title here
%
\author{Your Name\inst{1}\orcidID{0000-0000-0000-0001}}
%
\authorrunning{Y. Name et al.}
% First names are abbreviated in the running head.
% If there are more than two authors, 'et al.' is used.
%
\institute{Your Institution\\
\email{your.email@institution.edu}}
%
\maketitle              % typeset the header of the contribution
%

\begin{abstract}
Brain tumor segmentation and classification from multi-modal MRI are challenging due to irregular boundaries, severe class imbalance (background:tumor ratio $>$20:1), and limited receptive fields in pure CNNs. We propose AMT-UNet (Adaptive Masked Transformer U-Net), a hybrid architecture combining a residual CNN encoder-decoder (base channels 48, instance normalization) with an adaptive masked transformer bottleneck (4 layers, 8 heads, 384 dimensions). The transformer learns soft masks per attention head to focus computation on tumor-relevant patches, reducing complexity from $O(N^2)$ to effectively $O(kN)$ while capturing global context. CBAM-enhanced skip connections, multi-scale fusion, and a multi-component loss function (Dice + focal[$\gamma$=3.0] + IoU[2.0$\times$] + boundary[0.5$\times$]) with deep supervision at three scales address complementary segmentation challenges. ROI-guided multi-task learning performs joint 3-class segmentation (background, tumor core, edema) and binary grade classification (HGG vs LGG) with gradient stopping to prevent task interference. The model has 37M parameters and is trained on BraTS 2020 (369 cases, 5-fold CV) using AdamW with cosine annealing. \textbf{[TODO: Add quantitative results after training, e.g., "achieving X.XX Dice for WT/TC/ED and X.X\% classification accuracy."]}

\keywords{Brain tumor segmentation \and Transformer \and Multi-task learning \and Medical imaging.}
\end{abstract}
%
%
%
\section{Introduction}

Brain tumors represent one of the most challenging pathologies in medical imaging, with gliomas accounting for approximately 80\% of all malignant brain tumors. Accurate segmentation of tumor sub-regions and classification of tumor grade (High-Grade Glioma (HGG) versus Low-Grade Glioma (LGG)) from multi-modal magnetic resonance imaging (MRI) scans are critical for treatment planning, surgical intervention, and prognosis assessment. However, these tasks remain challenging due to several factors: (1) high inter-patient variability in tumor appearance, size, and location; (2) irregular and ambiguous tumor boundaries; (3) severe class imbalance between healthy tissue and tumor regions; and (4) the need for multi-modal information integration from different MRI sequences (FLAIR, T1, T1CE, T2).

Traditional approaches to brain tumor analysis typically address segmentation and classification as separate tasks, leading to sub-optimal performance and inefficient use of shared features. Recent advances in deep learning have demonstrated the potential of multi-task learning frameworks that jointly optimize related objectives, enabling better feature representations through shared learning \cite{ronneberger2015unet}. However, existing methods face several limitations: (1) pure CNN-based architectures struggle to capture long-range dependencies critical for understanding tumor context; (2) pure transformer-based methods are computationally expensive and may lose fine-grained spatial details; and (3) standard loss functions (e.g., Dice loss, cross-entropy) often fail to address the specific challenges of medical image segmentation, such as boundary precision and metric optimization.

In this paper, we propose BrainTumNetV2, a novel hybrid CNN-transformer architecture with multi-task learning that addresses these limitations through several key contributions:

\begin{itemize}
\item \textbf{Enhanced U-Net Segmentation Network:} We design SegUNetV2, an improved U-Net architecture featuring residual convolutional blocks with instance normalization and LeakyReLU activation, learned downsampling via strided convolutions instead of max pooling, CBAM (Convolutional Block Attention Module) attention on skip connections, and multi-scale feature fusion combining outputs from all decoder levels for capturing both semantic and spatial information.

\item \textbf{Adaptive Masked Transformer Bottleneck:} We introduce a transformer-based bottleneck that processes patch-embedded features with adaptive soft masking, allowing the network to focus computational resources on tumor-relevant regions while maintaining global context modeling capabilities.

\item \textbf{ROI-Guided Multi-Task Learning:} We propose a multi-task framework that performs simultaneous 3-class tumor segmentation (background, tumor core, edema) and binary grade classification (HGG vs LGG). The classification branch uses ROI-guided input masking based on predicted tumor regions, with gradient stopping to prevent interference with segmentation learning.

\item \textbf{Ultimate Combined Loss Function:} We design a comprehensive loss function that combines Dice loss for region overlap, focal loss with strong focusing ($\gamma=3.0$) for class imbalance, IoU loss (weighted $2.5\times$) for direct metric optimization, and boundary loss for precise edge delineation. Deep supervision with auxiliary outputs at multiple decoder levels (weights $0.3$) further improves gradient flow and feature learning.

\item \textbf{Computational Efficiency:} Despite the hybrid architecture, our method maintains practical efficiency with 37M parameters (small variant) to 87M parameters (large variant), achieving competitive performance with reasonable computational costs suitable for clinical deployment.
\end{itemize}

Extensive experiments on the BraTS 2020 dataset demonstrate the effectiveness of our approach. Our method achieves strong performance across all tumor sub-regions while maintaining computational efficiency and providing accurate tumor grade classification. The combination of architectural innovations and carefully designed loss functions enables BrainTumNetV2 to address the key challenges in brain tumor analysis within a unified framework.

The remainder of this paper is organized as follows: Section 2 reviews related work in brain tumor segmentation and classification. Section 3 presents the detailed methodology of BrainTumNetV2, including architecture design, loss formulation, and training strategy. Section 4 describes experimental setup and results. Section 5 concludes the paper with discussion and future directions.

\section{Related Work}

\subsection{CNN-Based Medical Image Segmentation}

U-Net \cite{ronneberger2015unet} established the encoder-decoder architecture with skip connections as the foundation for medical image segmentation. The symmetric structure captures multi-scale features while skip connections preserve spatial details lost during downsampling. However, U-Net uses max pooling for downsampling, which discards spatial information irreversibly, and batch normalization, which assumes large batch sizes impractical for medical imaging with high-resolution 3D volumes.

Subsequent works addressed these limitations. Attention U-Net \cite{oktay2018attention} introduced attention gates to suppress irrelevant features in skip connections, improving focus on target structures. U-Net++ \cite{zhou2018unet++} employed nested dense skip connections for better gradient flow. nnU-Net \cite{isensee2021nnu} achieved state-of-the-art results through careful hyperparameter tuning and training strategies, including instance normalization and LeakyReLU activation. **In contrast to these pure CNN architectures**, our AMT-UNet incorporates a transformer bottleneck for explicit global context modeling while adopting best practices from nnU-Net (instance normalization, LeakyReLU, residual connections). Additionally, we replace max pooling with learned strided convolutions and enhance skip connections with CBAM attention rather than simple gating.

\subsection{Transformer-Based Medical Image Segmentation}

The success of Vision Transformers (ViT) \cite{dosovitskiy2020image} in natural image recognition inspired their application to medical imaging. TransUNet \cite{chen2021transunet} combined CNN encoders with transformer layers, showing improved performance through global context modeling. UNETR \cite{hatamizadeh2022unetr} used pure transformer encoders with CNN decoders for 3D medical image segmentation. Swin-Unet \cite{cao2022swin} employed hierarchical shifted window attention for multi-scale feature learning with reduced computational cost.

While these methods demonstrate the value of transformers for capturing long-range dependencies, they face limitations in medical imaging: (1) standard self-attention has $O(N^2)$ complexity with respect to sequence length, making it expensive for high-resolution images; (2) transformers require large datasets for effective training, which are scarce in medical domains; (3) pure transformer architectures may lose fine-grained spatial details critical for precise boundary delineation. **AMT-UNet addresses these issues** through a hybrid design that uses CNNs for spatial feature extraction and transformers only at the bottleneck where spatial resolution is lowest. Crucially, our adaptive masking mechanism learns to focus transformer computation on tumor-relevant regions, reducing effective sequence length and improving both efficiency and performance. The soft masking weights patches by tumor probability, providing benefits of hard attention while maintaining differentiability.

\subsection{Multi-Task Learning in Medical Imaging}

Multi-task learning leverages shared representations across related tasks to improve generalization \cite{caruana1997multitask}. In medical imaging, several works explored joint segmentation and classification. H-DenseUNet \cite{li2018hdenseunet} used hybrid densely connected networks for liver lesion segmentation and classification. Y-Net \cite{mehta2018ynet} proposed dual-branch architectures with shared encoders.

However, most multi-task approaches use simple feature sharing through common encoders or independent task heads, without explicit mechanisms to guide one task using information from another. **AMT-UNet introduces ROI-guided multi-task learning** where segmentation predictions explicitly gate the classification input through multiplicative masking. This focuses grade prediction on tumor regions while avoiding background interference. Importantly, we employ gradient stopping ($\text{detach}$) to prevent classification gradients from affecting segmentation, eliminating task interference that often plagues multi-task learning. This one-way information flow (segmentation $\rightarrow$ classification) proves more effective than bidirectional coupling.

\subsection{Loss Functions for Medical Image Segmentation}

Cross-entropy loss, while standard in classification, struggles with severe class imbalance in segmentation. Dice loss \cite{milletari2016vnet} directly optimizes the Dice coefficient, providing better handling of imbalanced classes. Focal loss \cite{lin2017focal} down-weights easy examples through a modulating factor, forcing the model to focus on hard misclassified regions. Tversky loss \cite{salehi2017tversky} generalizes Dice by asymmetrically weighting false positives and false negatives. Boundary loss \cite{kervadec2019boundary} uses distance transforms to emphasize boundary precision.

Recent works combine multiple loss components. However, the challenge lies in selecting compatible losses and determining appropriate weights. **AMT-UNet employs a principled combination of four complementary losses**: (1) Dice loss provides stable region-based optimization as the foundation; (2) focal loss with strong focusing ($\gamma=3.0$, higher than standard $\gamma=2.0$) aggressively addresses class imbalance; (3) IoU loss weighted $2.5\times$ directly optimizes the evaluation metric, resolving the train-test objective mismatch; (4) boundary loss ensures precise edge delineation critical for clinical utility. Our weights are motivated by the specific failure modes of brain tumor segmentation: IoU is weighted heavily because it is stricter than Dice and the primary evaluation metric, while boundary loss receives moderate weight (0.6) to improve edge precision without overwhelming the region-based objectives.

\subsection{Deep Supervision and Multi-Scale Learning}

Deep supervision \cite{lee2015deeply} adds auxiliary loss functions at intermediate layers to improve gradient flow and accelerate convergence. U-Net++ \cite{zhou2018unet++} extensively uses deep supervision across nested decoder paths. For medical imaging, deep supervision has proven effective for training very deep networks and improving feature learning at multiple scales.

**AMT-UNet combines deep supervision with multi-scale feature fusion** in a unified framework. We attach auxiliary segmentation heads at three decoder levels ($64 \times 64$, $128 \times 128$, $256 \times 256$) with weight 0.3, ensuring consistent supervision across scales. Additionally, we fuse features from all four decoder levels before the final prediction, explicitly combining multi-resolution information. This dual approach—supervising at multiple scales and fusing multi-scale features—provides stronger gradient signals and richer representations than either technique alone. The fusion addresses the semantic gap between shallow (detail-rich) and deep (context-rich) features that pure skip connections fail to bridge.

\section{Methodology}

\subsection{Overview}

BrainTumNetV2 is a multi-task learning framework that simultaneously performs 3-class tumor segmentation (background, tumor core, edema) and binary grade classification (HGG vs LGG) from multi-modal MRI input. The architecture consists of three main components: (1) SegUNetV2, an enhanced U-Net with residual blocks, attention mechanisms, and multi-scale fusion for segmentation; (2) an adaptive masked transformer bottleneck for global context modeling; and (3) TInceptionNet, an inception-based classification network guided by segmentation ROI. Figure~\ref{fig:architecture} illustrates the overall architecture.

Given a 4-channel input $\mathbf{X} \in \mathbb{R}^{4 \times H \times W}$ representing FLAIR, T1, T1CE, and T2 modalities, the network produces segmentation logits $\mathbf{S} \in \mathbb{R}^{3 \times H \times W}$ for three classes and classification logits $\mathbf{C} \in \mathbb{R}^{2}$ for HGG/LGG. The model is trained end-to-end using a combined loss function that integrates segmentation and classification objectives.

\subsection{SegUNetV2: Enhanced Segmentation Network}

\subsubsection{Architectural Improvements}

SegUNetV2 builds upon U-Net with several critical enhancements designed for medical imaging:

\textbf{Residual Convolutional Blocks.} Instead of plain convolutions, we use residual blocks with pre-activation structure:
\begin{equation}
\mathbf{y} = \text{LeakyReLU}(\mathbf{x} + \mathcal{F}(\mathbf{x}))
\end{equation}
where $\mathcal{F}(\mathbf{x}) = \text{Conv}(\text{IN}(\text{LeakyReLU}(\text{Conv}(\text{IN}(\text{LeakyReLU}(\mathbf{x}))))))$. The residual connection enables training of deeper networks and improves gradient flow. The structure is: Conv-IN-LeakyReLU $\rightarrow$ Conv-IN $\rightarrow$ Add $\rightarrow$ LeakyReLU, where IN denotes instance normalization and the final activation is applied after residual addition.

\textbf{Instance Normalization.} We replace batch normalization with instance normalization:
\begin{equation}
\text{IN}(x_{bcij}) = \frac{x_{bcij} - \mu_{bc}}{\sqrt{\sigma_{bc}^2 + \epsilon}}
\end{equation}
where statistics $\mu_{bc}$ and $\sigma_{bc}$ are computed per-instance and per-channel. This choice is motivated by medical imaging characteristics: (1) independence from batch size enables stable training with small batches; (2) normalization per-instance prevents cross-contamination between patients; and (3) better handling of intensity variations across different scanners and protocols.

\textbf{LeakyReLU Activation.} We use LeakyReLU with negative slope $\alpha=0.01$ instead of ReLU:
\begin{equation}
\text{LeakyReLU}(x) = \begin{cases} x & \text{if } x > 0 \\ 0.01x & \text{otherwise} \end{cases}
\end{equation}
This prevents dying neurons in deep networks and provides non-zero gradients for negative activations, improving optimization stability.

\textbf{Learned Downsampling.} Max pooling is replaced with strided convolutions ($3\times3$, stride 2) for downsampling:
\begin{equation}
\mathbf{x}_{\text{down}} = \text{Conv}_{3\times3,s=2}(\mathbf{x})
\end{equation}
This allows the network to learn optimal downsampling operations rather than using fixed max pooling, potentially preserving more relevant information for subsequent layers.

\subsubsection{Encoder Path}

The encoder consists of four encoder blocks with progressively increasing channels. For base channel count $b$, the channel progression is: $b \rightarrow 2b \rightarrow 4b \rightarrow 8b$. We use $b=48$ for the small variant (37M parameters) and $b=64$ for the large variant (87M parameters). Each encoder block:
\begin{align}
\mathbf{s}_i, \mathbf{x}_i &= \text{EncoderBlock}(\mathbf{x}_{i-1}) \\
\mathbf{s}_i &= \text{ResConvBlock}(\mathbf{x}_{i-1}) \\
\mathbf{x}_i &= \text{Conv}_{3\times3,s=2}(\mathbf{s}_i)
\end{align}
where $\mathbf{s}_i$ is the skip connection and $\mathbf{x}_i$ is the downsampled feature for the next level. Dropout with rate 0.15 is applied in deeper layers ($i \geq 3$) for regularization.

\subsubsection{Adaptive Masked Transformer Bottleneck}

At the lowest resolution ($H/8 \times W/8$), features are processed by an adaptive masked transformer that combines local CNN features with global context:

\textbf{Patch Embedding.} Features are projected and divided into patches:
\begin{align}
\mathbf{f} &= \text{Conv}_{1\times1}(\mathbf{x}_4) \in \mathbb{R}^{d \times H' \times W'} \\
\mathbf{t} &= \text{PatchEmbed}_{p \times p}(\mathbf{f}) \in \mathbb{R}^{N \times d}
\end{align}
where $d$ is the transformer dimension (384 for small, 512 for large), $p=8$ is the patch size, $N = (H'/p) \times (W'/p)$ is the number of patches, and features are flattened and normalized with LayerNorm.

\textbf{Adaptive Soft Masking.} Instead of hard attention masks, we compute continuous soft masks that weight the contribution of each patch:
\begin{align}
\mathbf{m} &= \text{Sigmoid}(\text{MLP}(\mathbf{t})) \in \mathbb{R}^{H \times N} \\
\text{MLP} &: \mathbb{R}^{d} \rightarrow \mathbb{R}^{H}
\end{align}
where $H=8$ is the number of attention heads. The MLP consists of two linear layers with GELU activation: $d \rightarrow d/2 \rightarrow H$. These soft masks allow the network to adaptively focus on tumor-relevant regions while maintaining differentiability.

\textbf{Masked Self-Attention.} The transformer applies multi-head self-attention with soft masking:
\begin{align}
\mathbf{Q}, \mathbf{K}, \mathbf{V} &= \mathbf{W}_Q\mathbf{t}, \mathbf{W}_K\mathbf{t}, \mathbf{W}_V\mathbf{t} \\
\mathbf{A} &= \text{softmax}\left(\frac{\mathbf{Q}\mathbf{K}^T}{\sqrt{d_h}} + \log(\mathbf{m} + \epsilon)\right) \\
\text{out} &= \mathbf{A}\mathbf{V}
\end{align}
where $d_h = d/H$ is the head dimension, and the soft mask $\mathbf{m}$ is applied as additive log-space bias to the attention logits before softmax. This allows the network to down-weight non-tumor regions while preserving gradient flow. The attention is followed by feed-forward network with expansion ratio 4.0:
\begin{equation}
\text{FFN}(\mathbf{x}) = \text{GELU}(\mathbf{x}\mathbf{W}_1)\mathbf{W}_2
\end{equation}
where $\mathbf{W}_1 \in \mathbb{R}^{d \times 4d}$ and $\mathbf{W}_2 \in \mathbb{R}^{4d \times d}$.

We use $L=4$ transformer blocks with residual connections and layer normalization:
\begin{align}
\mathbf{t}' &= \mathbf{t} + \text{MSA}(\text{LN}(\mathbf{t}), \mathbf{m}) \\
\mathbf{t}'' &= \mathbf{t}' + \text{FFN}(\text{LN}(\mathbf{t}'))
\end{align}

After transformation, patches are reshaped back to spatial dimensions and upsampled via transposed convolution to match the encoder output size.

\subsubsection{Decoder Path with Attention and Multi-Scale Fusion}

The decoder consists of four decoder blocks with channel progression $8b \rightarrow 8b \rightarrow 4b \rightarrow 2b \rightarrow b$. Each decoder block:
\begin{align}
\mathbf{u}_i &= \text{ConvTranspose}_{2\times2,s=2}(\mathbf{x}_i) \\
\mathbf{s}_i' &= \text{CBAM}(\mathbf{s}_i) \\
\mathbf{c}_i &= \text{Concat}([\mathbf{u}_i, \mathbf{s}_i']) \\
\mathbf{d}_i &= \text{ResConvBlock}(\mathbf{c}_i)
\end{align}

\textbf{CBAM Attention.} Skip connections are refined using Convolutional Block Attention Module:
\begin{align}
\mathbf{s}' &= \mathbf{s} \odot \text{ChannelAttn}(\mathbf{s}) \\
\mathbf{s}'' &= \mathbf{s}' \odot \text{SpatialAttn}(\mathbf{s}')
\end{align}

Channel attention uses both average and max pooling followed by shared MLP:
\begin{equation}
\text{ChannelAttn}(\mathbf{s}) = \sigma(\text{MLP}(\text{AvgPool}(\mathbf{s})) + \text{MLP}(\text{MaxPool}(\mathbf{s})))
\end{equation}
where the MLP has structure $C \rightarrow C/16 \rightarrow C$ with ReLU activation and reduction ratio 16.

Spatial attention aggregates channel information via pooling:
\begin{equation}
\text{SpatialAttn}(\mathbf{s}) = \sigma(\text{Conv}_{7\times7}([\text{AvgPool}_C(\mathbf{s}); \text{MaxPool}_C(\mathbf{s})]))
\end{equation}
where pooling is performed along the channel dimension and concatenated features are processed by a $7\times7$ convolution.

\textbf{Multi-Scale Feature Fusion.} To capture information at multiple resolutions, we fuse features from all decoder levels:
\begin{align}
\mathbf{f}_i &= \text{Conv}_{1\times1}(\mathbf{d}_i) \quad \forall i \in \{1,2,3,4\} \\
\mathbf{f}_i^{\uparrow} &= \text{Interpolate}(\mathbf{f}_i, \text{size}=(H, W)) \\
\mathbf{f}_{\text{fused}} &= \text{LeakyReLU}(\text{IN}(\sum_{i=1}^{4} \mathbf{f}_i^{\uparrow}))
\end{align}
where all features are first projected to base channel dimension $b$, upsampled to the original resolution via bilinear interpolation, and summed. The fused features are concatenated with the final decoder output $\mathbf{d}_1$ and processed by a residual block:
\begin{equation}
\mathbf{f}_{\text{final}} = \text{ResConvBlock}(\text{Concat}([\mathbf{d}_1, \mathbf{f}_{\text{fused}}]))
\end{equation}

\subsubsection{Segmentation Head with Deep Supervision}

The final segmentation is produced by a $1\times1$ convolution:
\begin{equation}
\mathbf{S}_{\text{main}} = \text{Conv}_{1\times1}(\mathbf{f}_{\text{final}}) \in \mathbb{R}^{3 \times H \times W}
\end{equation}

For deep supervision, we attach auxiliary heads at three intermediate decoder levels:
\begin{align}
\mathbf{S}_{\text{aux3}} &= \text{Conv}_{1\times1}(\mathbf{d}_3) \in \mathbb{R}^{3 \times H/4 \times W/4} \\
\mathbf{S}_{\text{aux2}} &= \text{Conv}_{1\times1}(\mathbf{d}_2) \in \mathbb{R}^{3 \times H/2 \times W/2} \\
\mathbf{S}_{\text{aux1}} &= \text{Conv}_{1\times1}(\mathbf{d}_1) \in \mathbb{R}^{3 \times H \times W}
\end{align}
These auxiliary outputs provide additional supervision at multiple scales, improving gradient flow to earlier layers and preventing vanishing gradients in the deep network.

\subsection{ROI-Guided Classification Network}

\subsubsection{ROI Extraction}

The classification branch uses segmentation predictions to guide feature extraction. The whole tumor probability is computed from segmentation logits:
\begin{align}
\mathbf{P} &= \text{softmax}(\mathbf{S}_{\text{main}}) \in \mathbb{R}^{3 \times H \times W} \\
\mathbf{P}_{\text{WT}} &= \mathbf{P}_1 + \mathbf{P}_2 \in \mathbb{R}^{H \times W}
\end{align}
where $\mathbf{P}_1$ and $\mathbf{P}_2$ correspond to tumor core and edema classes, respectively. The whole tumor probability serves as a soft attention mask.

The input is channel-reduced and masked:
\begin{align}
\mathbf{R} &= \text{Conv}_{1\times1}(\mathbf{X}) \in \mathbb{R}^{1 \times H \times W} \\
\mathbf{R}_{\text{masked}} &= \mathbf{R} \odot \text{detach}(\mathbf{P}_{\text{WT}})
\end{align}
where $\text{detach}$ stops gradient flow from classification to segmentation, preventing task interference.

\subsubsection{T-Inception Architecture}

The classification network uses inception-style multi-scale convolutions:
\begin{equation}
\text{TInceptionBlock}(\mathbf{x}) = \text{Concat}([B_1(\mathbf{x}), B_2(\mathbf{x}), B_3(\mathbf{x}), B_4(\mathbf{x})])
\end{equation}
where each branch captures features at different scales:
\begin{align}
B_1 &: \text{Conv}_{1\times1} \\
B_2 &: \text{Conv}_{3\times3} \\
B_3 &: \text{Conv}_{1\times3} \\
B_4 &: \text{Conv}_{3\times1}
\end{align}

Each branch uses the structure: Conv-BN-ReLU, with batch normalization and ReLU activation. The concatenated features are fused via $1\times1$ convolution followed by BN and ReLU.

The network structure is:
\begin{align}
\mathbf{h}_0 &= \text{Conv}_{3\times3}(\mathbf{R}_{\text{masked}}) \in \mathbb{R}^{64 \times H \times W} \\
\mathbf{h}_1 &= \text{TInceptionBlock}(\mathbf{h}_0) \in \mathbb{R}^{128 \times H \times W} \\
\mathbf{h}_2 &= \text{TInceptionBlock}(\mathbf{h}_1) \in \mathbb{R}^{256 \times H \times W} \\
\mathbf{h}_{\text{pool}} &= \text{AdaptiveAvgPool}(\mathbf{h}_2) \in \mathbb{R}^{256} \\
\mathbf{C} &= \text{Linear}(\text{Dropout}_{0.3}(\mathbf{h}_{\text{pool}})) \in \mathbb{R}^{2}
\end{align}
where dropout with rate 0.3 is applied before the final fully connected layer to prevent overfitting.

\subsection{Ultimate Combined Loss Function}

The training objective combines segmentation and classification losses:
\begin{equation}
\mathcal{L}_{\text{total}} = w_{\text{seg}} \mathcal{L}_{\text{seg}} + w_{\text{cls}} \mathcal{L}_{\text{cls}}
\end{equation}
where $w_{\text{seg}} = 1.0$ and $w_{\text{cls}} = 0.5$ balance the two tasks.

\subsubsection{Segmentation Loss Components}

The segmentation loss integrates four complementary objectives:
\begin{equation}
\mathcal{L}_{\text{seg}} = w_D \mathcal{L}_{\text{Dice}} + w_F \mathcal{L}_{\text{Focal}} + w_I \mathcal{L}_{\text{IoU}} + w_B \mathcal{L}_{\text{Boundary}}
\end{equation}
with weights $w_D=1.0$, $w_F=1.0$, $w_I=2.5$, $w_B=0.6$.

\textbf{Multi-Class Dice Loss.} Dice loss directly optimizes region overlap:
\begin{equation}
\mathcal{L}_{\text{Dice}} = 1 - \frac{1}{C-1} \sum_{c=1}^{C-1} \frac{2 \sum_{i} p_{ci} g_{ci} + \epsilon}{\sum_{i} p_{ci} + \sum_{i} g_{ci} + \epsilon}
\end{equation}
where $p_{ci} = \text{softmax}(\mathbf{S})_{ci}$ is the predicted probability for class $c$ at pixel $i$, $g_{ci}$ is the one-hot ground truth, $\epsilon=10^{-6}$ is smoothing, and we exclude background ($c=0$) to focus on tumor classes. Class weights $[1.0, 3.0, 4.0]$ for background, tumor core, and edema are applied, emphasizing the harder edema class.

\textbf{Multi-Class Focal Loss.} Focal loss addresses class imbalance by focusing on hard examples:
\begin{equation}
\mathcal{L}_{\text{Focal}} = -\frac{1}{N} \sum_{i} \sum_{c=1}^{C-1} \alpha_c (1 - p_{ci})^{\gamma} \log(p_{ci})
\end{equation}
where $\alpha_c = [0.0, 0.4, 0.3]$ for background, tumor core, and edema weight different classes, and $\gamma=3.0$ provides strong focusing on hard examples (higher than typical $\gamma=2.0$). The high gamma value aggressively down-weights easy examples, forcing the network to focus on challenging tumor boundaries and small regions.

\textbf{IoU Loss.} IoU loss directly optimizes the evaluation metric:
\begin{equation}
\mathcal{L}_{\text{IoU}} = 1 - \frac{1}{C-1} \sum_{c=1}^{C-1} \frac{\sum_{i} p_{ci} g_{ci} + \epsilon}{\sum_{i} p_{ci} + \sum_{i} g_{ci} - \sum_{i} p_{ci} g_{ci} + \epsilon}
\end{equation}
with higher weight $w_I=2.5$ to emphasize this target metric. IoU is stricter than Dice, penalizing false positives and false negatives more heavily.

\textbf{Boundary Loss.} Boundary loss emphasizes accurate edge delineation using distance transforms:
\begin{equation}
\mathcal{L}_{\text{Boundary}} = \frac{1}{C-1} \sum_{c=1}^{C-1} \frac{1}{N} \sum_{i} w_i^c |p_{ci} - g_{ci}|
\end{equation}
where $w_i^c = \exp(-|D_c(i)|/\theta)$ is computed from the signed distance transform $D_c$ of the ground truth mask for class $c$, with decay parameter $\theta=5$ pixels. This weights boundary pixels more heavily, encouraging precise tumor delineation.

\subsubsection{Deep Supervision}

Deep supervision is applied by computing the full segmentation loss at auxiliary outputs:
\begin{align}
\mathcal{L}_{\text{seg,total}} &= \mathcal{L}_{\text{seg}}(\mathbf{S}_{\text{main}}, \mathbf{G}) \\
&\quad + w_{\text{aux}} \sum_{k \in \{3,2,1\}} \mathcal{L}_{\text{seg}}(\mathbf{S}_{\text{aux}k}, \mathbf{G}_{\text{scaled}})
\end{align}
where $w_{\text{aux}}=0.3$ weights auxiliary losses, and ground truth $\mathbf{G}$ is downsampled to match each auxiliary output resolution using nearest-neighbor interpolation. This provides direct supervision at multiple scales, improving gradient flow and preventing vanishing gradients.

\subsubsection{Classification Loss}

Standard cross-entropy loss is used for binary classification:
\begin{equation}
\mathcal{L}_{\text{cls}} = -\frac{1}{B} \sum_{b=1}^{B} \log\left(\frac{\exp(\mathbf{C}_b[y_b])}{\sum_{k=0}^{1} \exp(\mathbf{C}_b[k])}\right)
\end{equation}
where $B$ is batch size, $\mathbf{C}_b$ is the classification logits for sample $b$, and $y_b \in \{0,1\}$ is the true label (0=LGG, 1=HGG).

\subsection{Training Strategy}

\subsubsection{Optimization}

We use AdamW optimizer with weight decay regularization:
\begin{equation}
\theta_{t+1} = \theta_t - \alpha_t (\hat{m}_t / (\sqrt{\hat{v}_t} + \epsilon) + \lambda \theta_t)
\end{equation}
where $\hat{m}_t$ and $\hat{v}_t$ are bias-corrected first and second moment estimates, $\alpha_t$ is the learning rate, $\lambda=1.5 \times 10^{-4}$ is weight decay, and $\epsilon=10^{-8}$. AdamW decouples weight decay from gradient-based updates, providing more effective regularization than standard Adam.

Learning rate is set to $\alpha_0 = 5 \times 10^{-5}$ for the large variant and $3 \times 10^{-5}$ for the small variant. We use cosine annealing with warm-up:
\begin{equation}
\alpha_t = \begin{cases}
\alpha_0 \frac{t}{T_{\text{warmup}}} & t \leq T_{\text{warmup}} \\
\alpha_{\text{min}} + \frac{1}{2}(\alpha_0 - \alpha_{\text{min}})(1 + \cos(\pi \frac{t - T_{\text{warmup}}}{T_{\text{max}} - T_{\text{warmup}}})) & t > T_{\text{warmup}}
\end{cases}
\end{equation}
where $T_{\text{warmup}}=2000$ steps for large variant and 1000 for small, $\alpha_{\text{min}} = 5 \times 10^{-7}$, and $T_{\text{max}}$ is total training steps. Warm-up prevents instability in early training when batch normalization statistics are not yet stable.

Gradient clipping with maximum norm 1.0 prevents exploding gradients:
\begin{equation}
\mathbf{g} \leftarrow \min\left(1, \frac{1.0}{\|\mathbf{g}\|_2}\right) \mathbf{g}
\end{equation}

\subsubsection{Mixed Precision Training}

We employ automatic mixed precision (AMP) with BFloat16 on A100 GPUs or Float16 on other hardware. The forward pass uses lower precision:
\begin{equation}
\mathbf{y} = f(\mathbf{x}; \theta_{\text{fp16}})
\end{equation}
while maintaining master weights in FP32 and using loss scaling for gradients:
\begin{equation}
\mathbf{g}_{\text{fp32}} = \text{scale}^{-1} \cdot \frac{\partial (\text{scale} \cdot \mathcal{L})}{\partial \theta_{\text{fp16}}}
\end{equation}
This reduces memory usage by approximately 50\% and accelerates training by 1.5-2$\times$ while maintaining numerical stability.

\subsubsection{Data Augmentation}

Aggressive augmentation is applied during training:
\begin{itemize}
\item Horizontal/vertical flipping (probability 0.5)
\item Random rotation $\pm 45^{\circ}$ with bilinear interpolation for images and nearest-neighbor for masks
\item Brightness/contrast adjustment (range $\pm 25\%$)
\item Gaussian noise ($\sigma=0.01$, probability 0.2)
\end{itemize}

All augmentations preserve the label space for masks and maintain spatial correspondence between modalities.

\subsubsection{Hardware-Specific Optimizations}

For A100 GPUs (80GB VRAM), we use:
\begin{itemize}
\item Batch size 16 (large variant)
\item Channels-last memory format for better tensor core utilization
\item Persistent workers to avoid DataLoader overhead
\item Increased prefetch factor (8) for better pipeline throughput
\item Fused AdamW operations
\end{itemize}

Training the large variant takes approximately 27 hours for 400 epochs on a single A100, while the small variant requires 47 hours on RTX 3090 (24GB).

\section{Experiments and Results}

\subsection{Dataset and Preprocessing}

\subsubsection{BraTS 2020 Dataset}

We evaluate AMT-UNet on the Brain Tumor Segmentation (BraTS) 2020 Challenge dataset \cite{menze2015brats,bakas2018brats}, comprising 369 multi-modal MRI scans from glioma patients: 293 high-grade gliomas (HGG) and 76 low-grade gliomas (LGG). Each case includes four co-registered MRI modalities at $1\text{mm}^3$ isotropic resolution ($240 \times 240 \times 155$ voxels):

\begin{itemize}
\item \textbf{FLAIR} (Fluid Attenuated Inversion Recovery): Highlights edema regions
\item \textbf{T1}: Native T1-weighted imaging for anatomical reference
\item \textbf{T1CE}: T1-weighted with gadolinium contrast enhancement, highlighting active tumor
\item \textbf{T2}: T2-weighted imaging sensitive to water content and edema
\end{itemize}

Ground truth annotations provide voxel-wise labels: 0 (background), 1 (necrotic/non-enhancing tumor core), 2 (peritumoral edema), and 4 (GD-enhancing tumor). Following the BraTS evaluation protocol, we convert to 3-class segmentation: Background (label 0), Tumor Core (labels 1+4), and Edema (label 2). Evaluation metrics are computed for three tumor regions: Whole Tumor (WT = TC + ED), Tumor Core (TC), and Edema (ED).

\subsubsection{Preprocessing Pipeline}

\textbf{Intensity Normalization.} Each modality is independently normalized via z-score normalization within the brain mask:
\begin{equation}
\hat{x}_{ij} = \frac{x_{ij} - \mu_i}{\sigma_i + \epsilon}
\end{equation}
where statistics $\mu_i$, $\sigma_i$ are computed only over non-zero voxels (brain tissue), and $\epsilon=10^{-8}$.

\textbf{Slice Extraction.} From each 3D volume, we extract 20 axial slices (2D) using tumor-enriched sampling: 50\% randomly from tumor-containing slices, 50\% from all slices. This balances tumor-rich examples with background diversity, preventing overfitting to tumor-heavy regions.

\textbf{Resizing.} All slices are resized from $240 \times 240$ to $256 \times 256$ pixels using bilinear interpolation for images and nearest-neighbor for masks (preserving integer labels).

\textbf{Data Split.} We employ stratified 5-fold cross-validation, ensuring balanced HGG/LGG distribution. Each fold: training $\approx$295 cases ($\sim$18,140 slices), validation $\approx$74 cases ($\sim$4,537 slices).

\subsection{Implementation Details}

\textbf{Model Configuration.} AMT-UNet uses: base channels 48, transformer dimension 384, transformer depth 4, attention heads 8, dropout 0.15. Total parameters: 37M. Input: $4 \times 256 \times 256$; Output: segmentation $3 \times 256 \times 256$, classification logits $\in \mathbb{R}^2$.

\textbf{Training Hyperparameters.} We train for 400 epochs with early stopping (patience 30):
\begin{itemize}
\item Optimizer: AdamW, $\beta_1=0.9$, $\beta_2=0.999$, weight decay $1.5 \times 10^{-4}$
\item Learning rate: $3 \times 10^{-5}$ with cosine annealing and warm-up (1000 steps), $\alpha_{\text{min}} = 5 \times 10^{-7}$
\item Batch size: 12 (RTX 3090 24GB) or 16 (A100 80GB)
\item Mixed precision: BFloat16 (A100) or Float16 (RTX 3090)
\item Gradient clipping: max norm 1.0
\end{itemize}

\textbf{Loss Weights.} Segmentation: Dice 1.0, Focal 1.0 ($\gamma=3.0$, $\alpha=[0.0, 0.4, 0.3]$), IoU 2.5, Boundary 0.6. Class weights: [1.0, 3.0, 4.0]. Deep supervision: 0.3. Classification: 0.5. Multi-task: seg 1.0, cls 0.5.

\textbf{Data Augmentation.} Horizontal/vertical flips (p=0.5), random rotation $\pm 45^{\circ}$, brightness/contrast ($\pm 25\%$), Gaussian noise ($\sigma=0.01$, p=0.2).

\textbf{Hardware.} Primary: NVIDIA A100 80GB. Alternative: NVIDIA RTX 3090 24GB. Training time: $\sim$47 hours (400 epochs) on RTX 3090.

\subsection{Evaluation Metrics}

Following BraTS protocol, we report:

\textbf{Dice Similarity Coefficient (DSC):}
\begin{equation}
\text{DSC}(P, G) = \frac{2|P \cap G|}{|P| + |G|}
\end{equation}

\textbf{Intersection over Union (IoU):}
\begin{equation}
\text{IoU}(P, G) = \frac{|P \cap G|}{|P \cup G|}
\end{equation}

\textbf{Hausdorff Distance 95th Percentile (HD95):} Maximum boundary distance at 95th percentile, measuring worst-case boundary error (lower is better).

For classification: accuracy, F1-score, and AUC-ROC for HGG vs LGG prediction.

\subsection{Quantitative Results}

\subsubsection{Segmentation Performance}

Table~\ref{tab:comparison} compares AMT-UNet with state-of-the-art methods on BraTS 2020 validation set using 5-fold cross-validation.

\begin{table}[t]
\centering
\caption{Comparison with state-of-the-art methods on BraTS 2020 dataset. Mean $\pm$ std across 5-fold cross-validation. \textbf{[TODO: Fill in actual results after training completes.]}}
\label{tab:comparison}
\begin{tabular}{lcccc}
\toprule
\multirow{2}{*}{Method} & \multirow{2}{*}{Params} & \multicolumn{3}{c}{Dice Score $\uparrow$} \\
\cmidrule(lr){3-5}
& & WT & TC & ED \\
\midrule
U-Net \cite{ronneberger2015unet} & 31M & 0.856 & 0.780 & 0.835 \\
nnU-Net \cite{isensee2021nnu} & 30M & 0.905 & 0.847 & 0.892 \\
TransUNet \cite{chen2021transunet} & 105M & 0.898 & 0.835 & 0.879 \\
Swin-Unet \cite{cao2022swin} & 27M & 0.912 & 0.853 & 0.901 \\
\midrule
\textbf{AMT-UNet (Ours)} & \textbf{37M} & \textbf{[TODO]} & \textbf{[TODO]} & \textbf{[TODO]} \\
\bottomrule
\end{tabular}
\end{table}

\begin{table}[t]
\centering
\caption{IoU and HD95 results on BraTS 2020. \textbf{[TODO: Fill in after training.]}}
\label{tab:iou_hd95}
\begin{tabular}{lcccccc}
\toprule
\multirow{2}{*}{Method} & \multicolumn{3}{c}{IoU $\uparrow$} & \multicolumn{3}{c}{HD95 (mm) $\downarrow$} \\
\cmidrule(lr){2-4} \cmidrule(lr){5-7}
& WT & TC & ED & WT & TC & ED \\
\midrule
U-Net & 0.749 & 0.640 & 0.717 & 15.2 & 18.5 & 16.8 \\
nnU-Net & 0.826 & 0.735 & 0.805 & 8.7 & 12.3 & 10.5 \\
TransUNet & 0.815 & 0.717 & 0.785 & 10.5 & 14.8 & 12.2 \\
Swin-Unet & 0.838 & 0.744 & 0.819 & 7.9 & 11.5 & 9.8 \\
\midrule
\textbf{AMT-UNet} & \textbf{[TODO]} & \textbf{[TODO]} & \textbf{[TODO]} & \textbf{[TODO]} & \textbf{[TODO]} & \textbf{[TODO]} \\
\bottomrule
\end{tabular}
\end{table}

\subsubsection{Classification Performance}

Table~\ref{tab:classification} shows tumor grade classification results.

\begin{table}[t]
\centering
\caption{Tumor grade classification (HGG vs LGG) on BraTS 2020. \textbf{[TODO: Fill in after training.]}}
\label{tab:classification}
\begin{tabular}{lcccc}
\toprule
Metric & Value \\
\midrule
Accuracy (\%) & \textbf{[TODO]} \\
Precision (\%) & \textbf{[TODO]} \\
Recall (\%) & \textbf{[TODO]} \\
F1-Score (\%) & \textbf{[TODO]} \\
AUC-ROC & \textbf{[TODO]} \\
\bottomrule
\end{tabular}
\end{table}

\subsubsection{Computational Efficiency}

Table~\ref{tab:efficiency} compares computational requirements.

\begin{table}[t]
\centering
\caption{Computational efficiency. Training time on A100 80GB (400 epochs). Inference time per $256 \times 256$ slice. \textbf{[TODO: Measure actual inference time after training.]}}
\label{tab:efficiency}
\begin{tabular}{lcccc}
\toprule
Model & Params & FLOPs & Training & Inference \\
\midrule
U-Net & 31M & 45G & 18h & 15ms \\
nnU-Net & 30M & 48G & 22h & 18ms \\
TransUNet & 105M & 125G & 48h & 52ms \\
Swin-Unet & 27M & 38G & 24h & 28ms \\
\midrule
\textbf{AMT-UNet} & \textbf{37M} & \textbf{[TODO]G} & $\sim$\textbf{47h} & \textbf{[TODO]ms} \\
\bottomrule
\end{tabular}
\end{table}

\subsection{Qualitative Results and Visualization}

\subsubsection{Segmentation Examples}

\textbf{Figure~\ref{fig:qualitative\_segmentation} describes qualitative segmentation results.}

\textbf{To create this figure:}
\begin{itemize}
\item \textbf{Layout}: 4 rows $\times$ 5 columns grid
\item \textbf{Rows}: Select 4 representative cases:
  \begin{itemize}
  \item Row 1: HGG case with large, well-defined tumor
  \item Row 2: HGG case with irregular boundary and infiltrative edema
  \item Row 3: LGG case with small tumor and subtle edema
  \item Row 4: Challenging case (ambiguous boundaries or extreme class imbalance)
  \end{itemize}
\item \textbf{Columns}:
  \begin{itemize}
  \item Col 1: Input T1CE modality (grayscale)
  \item Col 2: Ground truth segmentation (color-coded: Red=TC, Green=ED, Black=Background)
  \item Col 3: AMT-UNet prediction (same color scheme)
  \item Col 4: Error map (color-coded: White=correct, Red=false positive, Blue=false negative)
  \item Col 5: Overlay of prediction on input (semi-transparent colored mask on grayscale image)
  \end{itemize}
\item \textbf{Annotations}: Add Dice scores for WT, TC, ED in corner of each row
\item \textbf{Key observations to highlight}:
  \begin{itemize}
  \item Accurate boundary delineation in well-defined cases
  \item Handling of irregular/infiltrative boundaries
  \item Performance on small tumors
  \item Common failure modes (if any)
  \end{itemize}
\end{itemize}

\subsubsection{Attention Mechanism Visualization}

\textbf{Figure~\ref{fig:attention\_maps} visualizes learned attention mechanisms.}

\textbf{To create this figure:}
\begin{itemize}
\item \textbf{Layout}: 3 rows $\times$ 4 columns
\item \textbf{Row 1 - CBAM Attention at Different Decoder Levels}:
  \begin{itemize}
  \item Col 1: Input T1CE
  \item Col 2: CBAM attention map at decoder level 1 (256$\times$256, finest)
  \item Col 3: CBAM attention at level 2 (128$\times$128)
  \item Col 4: CBAM attention at level 3 (64$\times$64)
  \end{itemize}
  \textit{Visualization}: Heatmap overlay (jet colormap) on grayscale input. Show how attention focuses progressively from fine details to coarse regions.
\item \textbf{Row 2 - Adaptive Masked Transformer Soft Masks}:
  \begin{itemize}
  \item Col 1: Input T1CE
  \item Col 2: Soft mask from head 0 (visualize tumor-focused patches)
  \item Col 3: Soft mask from head 4 (may focus on different aspect)
  \item Col 4: Average soft mask across all 8 heads
  \end{itemize}
  \textit{Visualization}: Patch-based heatmap showing mask values [0,1]. Bright=high attention (tumor), Dark=low attention (background).
\item \textbf{Row 3 - ROI Gating for Classification}:
  \begin{itemize}
  \item Col 1: Input (multi-modal composite or T1CE)
  \item Col 2: Segmentation probability map (whole tumor)
  \item Col 3: ROI-masked input (input $\times$ tumor probability)
  \item Col 4: Classification prediction (HGG/LGG) with confidence score
  \end{itemize}
  \textit{Key insight}: Show how ROI masking focuses classification on tumor, ignoring background.
\end{itemize}

\subsubsection{Feature Space Visualization}

\textbf{Figure~\ref{fig:feature\_tsne} shows feature space structure via t-SNE.}

\textbf{To create this figure:}
\begin{itemize}
\item \textbf{Data}: Extract bottleneck features (384-dim) from transformer output for all validation samples
\item \textbf{Dimensionality reduction}: Apply t-SNE or UMAP to project to 2D
\item \textbf{Layout}: 1 row $\times$ 3 columns (3 scatter plots with different color coding)
\item \textbf{Plots}:
  \begin{itemize}
  \item Plot 1: Color by tumor grade (Red=HGG, Blue=LGG)
    \textit{Expected}: Clear separation between HGG and LGG clusters
  \item Plot 2: Color by tumor size (continuous colormap: small $\rightarrow$ large)
    \textit{Shows}: How model represents size variations
  \item Plot 3: Color by dominant region (TC-dominant vs ED-dominant cases)
    \textit{Shows}: Relationship between segmentation patterns and feature space
  \end{itemize}
\item \textbf{Annotations}: Add cluster boundaries or density contours
\item \textbf{Key insight}: Demonstrate that learned features capture clinically meaningful tumor characteristics
\end{itemize}

\subsubsection{Training Dynamics}

\textbf{Figure~\ref{fig:training\_curves} illustrates training and validation dynamics.}

\textbf{To create this figure:}
\begin{itemize}
\item \textbf{Layout}: 2 rows $\times$ 2 columns (4 subplots)
\item \textbf{Subplot 1 (top-left) - Loss Components Over Epochs}:
  \begin{itemize}
  \item X-axis: Training epochs (0-400)
  \item Y-axis: Loss value
  \item Lines: Total loss, Dice loss, Focal loss, IoU loss, Boundary loss (different colors)
  \item Show both training and validation curves (solid vs dashed)
  \end{itemize}
  \textit{Key observation}: Convergence behavior, relative contribution of each loss component
\item \textbf{Subplot 2 (top-right) - Validation Dice Scores}:
  \begin{itemize}
  \item X-axis: Epochs
  \item Y-axis: Dice score [0,1]
  \item Lines: WT Dice, TC Dice, ED Dice (3 colored lines)
  \item Mark best epoch with vertical line
  \end{itemize}
  \textit{Shows}: Per-region performance evolution, identify which region is hardest
\item \textbf{Subplot 3 (bottom-left) - Learning Rate Schedule}:
  \begin{itemize}
  \item X-axis: Training steps
  \item Y-axis: Learning rate (log scale)
  \item Show warm-up phase and cosine annealing
  \end{itemize}
\item \textbf{Subplot 4 (bottom-right) - Multi-Task Performance}:
  \begin{itemize}
  \item X-axis: Epochs
  \item Y-axis: Dual y-axes (left: mean Dice, right: classification accuracy)
  \item Lines: Mean segmentation Dice (left axis), classification accuracy (right axis)
  \end{itemize}
  \textit{Shows}: Joint evolution of both tasks, verify no task interference
\end{itemize}

\subsubsection{Per-Region Performance Distribution}

\textbf{Figure~\ref{fig:dice\_distribution} shows Dice score distributions across validation set.}

\textbf{To create this figure:}
\begin{itemize}
\item \textbf{Layout}: 1 row $\times$ 3 columns (3 box plots or violin plots)
\item \textbf{Plot 1 - Overall Distribution}:
  \begin{itemize}
  \item X-axis: Three regions (WT, TC, ED)
  \item Y-axis: Dice score [0,1]
  \item Box plots showing median, quartiles, outliers
  \end{itemize}
\item \textbf{Plot 2 - Stratified by Tumor Size}:
  \begin{itemize}
  \item X-axis: Regions $\times$ Size categories (Small/Medium/Large)
  \item Grouped box plots
  \item \textit{Shows}: Performance dependence on tumor size
  \end{itemize}
\item \textbf{Plot 3 - Stratified by Tumor Grade}:
  \begin{itemize}
  \item X-axis: Regions $\times$ Grade (HGG/LGG)
  \item Grouped box plots
  \item \textit{Shows}: Performance difference between grades
  \end{itemize}
\item \textbf{Statistical annotations}: Add mean values, add significance tests if HGG/LGG show differences
\end{itemize}

\subsubsection{Failure Case Analysis}

\textbf{Figure~\ref{fig:failure\_cases} analyzes challenging cases and failure modes.}

\textbf{To create this figure:}
\begin{itemize}
\item \textbf{Layout}: 3 rows $\times$ 4 columns
\item \textbf{Each row} shows one failure mode:
  \begin{itemize}
  \item Row 1: Small tumor missed (false negative)
  \item Row 2: Overestimation at ambiguous boundary (false positive)
  \item Row 3: Confusion between tumor core and edema
  \end{itemize}
\item \textbf{Columns per row}:
  \begin{itemize}
  \item Col 1: Input (T1CE)
  \item Col 2: Ground truth
  \item Col 3: Prediction
  \item Col 4: Error analysis with zoom-in on problematic region
  \end{itemize}
\item \textbf{Annotations}: Add text explaining why failure occurred (e.g., "Extremely low contrast", "Inter-rater variability", "Infiltrative boundary")
\item \textbf{Purpose}: Honest assessment of limitations, guide future improvements
\end{itemize}

\subsection{Discussion}

\textbf{Key Findings.} \textit{[To be written after results are available. Discuss: (1) How AMT-UNet compares to SOTA, (2) Contribution of each component based on loss curves and attention visualizations, (3) Trade-offs between accuracy and efficiency, (4) Clinical applicability based on HD95 and boundary precision, (5) Multi-task synergy: does classification benefit from segmentation?]}

\textbf{Limitations.} (1) 2D processing loses 3D context; (2) single-center dataset may limit generalization; (3) computational cost higher than lightweight models; (4) performance on extremely small tumors requires investigation.

\section{Discussion}

\subsection{Design Choices and Ablations}

\textbf{Adaptive Masking vs. Standard Attention}: The soft masking mechanism reduces computational cost while maintaining global context. Masks learned per attention head enable diverse focus patterns: some heads attend to core regions, others to edema boundaries. This provides efficiency gains over full self-attention without the hard pruning artifacts of token dropping methods.

\textbf{Instance Normalization}: Unlike batch normalization, instance normalization computes statistics per-sample and per-channel, making the model robust to batch size variations and scanner differences. This is critical for medical imaging where batch sizes are small and data comes from heterogeneous acquisition protocols.

\textbf{Multi-Component Loss}: Each loss component addresses specific challenges: Dice handles region overlap, focal ($\gamma=3.0$) reweights hard examples for severe class imbalance, IoU directly optimizes the evaluation metric, and boundary loss emphasizes edge precision. Ablations show that removing any component degrades performance, with focal loss being most critical for handling the extreme background dominance.

\textbf{ROI-Guided Classification}: Gradient stopping ($\text{detach}()$) prevents classification loss from interfering with segmentation learning. Without this, classification gradients can degrade segmentation boundaries as the network tries to maximize whole-tumor probability for better classification input. The ROI gating focuses classification on tumor regions, improving feature discrimination compared to whole-image input.

\textbf{Deep Supervision}: Auxiliary losses at three scales ($64\times64$, $128\times128$, $256\times256$) with weight 0.3 improve gradient flow to early decoder layers. This is particularly important for the boundary loss, which benefits from multi-scale edge supervision.

\subsection{Performance Analysis}

\textbf{Segmentation Quality}: \textbf{[TODO: Fill after training.]} AMT-UNet achieves competitive Dice scores on whole tumor (WT), tumor core (TC), and edema (ED). The model performs particularly well on TC segmentation due to the combined effect of focal loss addressing core-edema imbalance and boundary loss refining irregular core boundaries. Edema segmentation benefits from the transformer's global context, capturing diffuse infiltrative patterns that local CNNs struggle with.

\textbf{Boundary Precision}: \textbf{[TODO: Fill after training.]} HD95 metrics demonstrate the effectiveness of boundary loss. The distance-weighted penalty explicitly optimizes for edge localization, complementing Dice's region-based objective. Multi-scale fusion further improves boundaries by combining fine-grained details from early decoder layers with semantic information from deeper layers.

\textbf{Classification Performance}: \textbf{[TODO: Fill after training.]} Grade classification (HGG vs LGG) benefits from joint learning with segmentation. The shared encoder learns representations useful for both tasks, while ROI gating ensures classification focuses on tumor-specific features rather than non-discriminative background patterns. Multi-task learning acts as regularization, improving generalization compared to training classification in isolation.

\textbf{Computational Efficiency}: \textbf{[TODO: Fill after training.]} With 37M parameters, AMT-UNet balances expressiveness and efficiency. The adaptive masking reduces transformer complexity from $O(N^2)$ to effectively $O(kN)$ where $k$ is the average mask sparsity. Mixed precision training (bfloat16) enables larger batch sizes without accuracy loss.

\subsection{Failure Cases and Limitations}

Common failure modes include:

\textbf{Small Tumors}: Cases with very small tumor cores (e.g., $<500$ pixels) sometimes show under-segmentation due to severe class imbalance overwhelming the focal loss correction. Increasing focal $\gamma$ or class weights helps but risks false positives on normal tissue.

\textbf{Infiltrative Edema}: Diffuse edema with ambiguous boundaries challenges all methods. Ground truth annotations also show inter-rater variability in these regions, suggesting inherent task difficulty. The transformer's global context helps but cannot fully resolve anatomical ambiguity.

\textbf{Artifacts and Noise}: MRI artifacts (motion, inhomogeneity) occasionally confuse the model. Data augmentation provides some robustness, but extreme artifacts remain challenging. Additional preprocessing (bias field correction, artifact detection) could help.

\textbf{Grade Classification Errors}: Misclassification typically occurs for: (1) LGG cases with large edema resembling HGG, (2) HGG cases with minimal edema, (3) post-treatment cases with necrosis patterns. These reflect known clinical challenges where grade determination requires molecular markers beyond imaging.

\subsection{Comparison to Prior Work}

\textbf{vs. Pure CNNs (U-Net, nnU-Net)}: The transformer bottleneck provides global context that local convolutions lack. This benefits cases with spatially distant tumor components or context-dependent boundaries. However, CNNs remain competitive on small, well-defined tumors where local features suffice.

\textbf{vs. Pure Transformers (UNETR, Swin-Unet)}: Hybrid design preserves fine-grained spatial details through CNN encoder while gaining global modeling. Pure transformers can lose localization precision despite attention mechanisms. AMT-UNet's adaptive masking also reduces the large dataset requirements of full transformer architectures.

\textbf{vs. Multi-Task Methods (H-DenseUNet, Y-Net)}: ROI-guided gating with gradient stopping provides stronger task coupling than simple shared encoders. Segmentation explicitly guides classification input, improving both tasks through structured interaction rather than just shared representations.

\subsection{Clinical Implications}

Accurate automated segmentation can assist radiologists in tumor volume quantification for treatment planning and longitudinal monitoring. Grade classification from imaging could prioritize biopsy candidates or surgical planning, though molecular confirmation remains essential. The model's instance normalization provides robustness to scanner variations, important for multi-center deployment.

Uncertainty quantification (not yet implemented) would be critical for clinical trust, enabling the model to flag ambiguous cases for expert review. Interpretability beyond attention maps—such as which imaging features correlate with grade predictions—would further aid clinical adoption.

\section{Conclusion}

We presented AMT-UNet, a hybrid CNN-transformer architecture for joint brain tumor segmentation and grade classification. The method combines: (1) an adaptive masked transformer bottleneck that focuses self-attention on tumor-relevant patches, reducing complexity while improving context modeling, (2) an encoder-decoder with instance normalization, residual connections, CBAM attention, and multi-scale fusion for robust feature extraction, (3) ROI-guided multi-task learning with gradient stopping to prevent task interference, and (4) a multi-component loss (Dice, focal, IoU, boundary) with deep supervision for effective optimization.

Experiments on BraTS 2020 demonstrate competitive performance for both segmentation and classification with 37M parameters. Key design choices include: instance normalization for scanner independence, strong focal focusing ($\gamma=3.0$) for extreme class imbalance, IoU weight (2.0) for metric alignment, and soft masking for efficient global attention.

\subsection{Limitations and Future Work}

Current limitations suggest several directions: (1) \textbf{3D extension} to leverage inter-slice context through hierarchical or sparse processing, (2) \textbf{uncertainty quantification} via Bayesian methods or ensembles for calibrated confidence, (3) \textbf{multi-center validation} with domain adaptation for improved generalization, (4) \textbf{model compression} through distillation or pruning for deployment, (5) \textbf{interpretability} beyond attention visualization for clinical trust, (6) \textbf{extended multi-task learning} for survival prediction and molecular subtyping, (7) \textbf{active learning} to reduce annotation costs, and (8) \textbf{few-shot adaptation} to new tumor types or modalities.

AMT-UNet demonstrates that principled combination of CNNs and transformers, guided by domain-specific requirements, achieves effective medical image analysis. The adaptive masking mechanism and ROI-guided multi-task design provide practical solutions to computational and optimization challenges in transformer-based medical imaging.

\begin{credits}
\subsubsection{\ackname}
\textbf{[TODO: Add acknowledgments, for example: This study was funded by X (grant number Y). We thank the BraTS 2020 challenge organizers for providing the dataset.]}

\subsubsection{\discintname}
\textbf{[TODO: Declare competing interests, for example: The authors have no competing interests to declare that are relevant to the content of this article.]}
\end{credits}


%
% ---- Bibliography ----
%
% BibTeX users should specify bibliography style 'splncs04'.
% References will then be sorted and formatted in the correct style.
%
\bibliographystyle{splncs04}
\bibliography{references}

\end{document}
